% STATUS: written 2026-11-06 (English gate clean; sources read in full as recorded below).
% Sources read in full: Ore 1933 (Ann. of Math. 34, 480--508).
\section[Skew polynomials and twisted derivations: Ore 1933][Skew polynomials and twisted derivations]{Skew polynomials and twisted derivations: Ore 1933}\label{sec:origins:ore-and-twists}

The next step came from polynomials. In the paper \emph{Theory of non-commutative polynomials}~\cite{Ore1933}, Ore
builds one theory for polynomials over a field $K$. The field may be
commutative or not. There is only one condition. The degree of a product must
be the sum of the degrees of the factors. Ore says that this one theory
contains ``not only the equations studied in non-commutative algebras but
also most of the linear operational equations of analysis as, for instance,
linear differential and difference equations''~\cite[p.~480]{Ore1933}.

He notes that the theory could also come from the work of Noether and
Schmeidler. They studied modules over rings of differential and difference
expressions. Ore instead builds the theory directly from the
polynomials~\cite[p.~480]{Ore1933}.

How should such polynomials multiply? The degree condition fixes the shape of the multiplication. It is enough to
define $x\cdot a$ for $a \in K$. The condition leaves only one possibility.

\begin{definition}[{\cite[p.~481, (3)]{Ore1933}}]\label{def:origins:ore}
For $a \in K$,
\[
x\cdot a = \bar a\, x + a',
\]
where $\bar a$ and $a'$ are elements of $K$. Ore calls $\bar a$ the
\emph{conjugate} and $a'$ the \emph{derivative} of $a$.
\end{definition}

The laws of associativity and distributivity give four rules:
\[
\overline{a+b} = \bar a + \bar b,\quad \overline{ab} = \bar a\,\bar b,\qquad
(a+b)' = a' + b',\quad (ab)' = \bar a\, b' + a' b
\]
\cite[p.~482, (4)--(5)]{Ore1933}. So the conjugation is an endomorphism of
$K$. The derivative satisfies a product rule twisted by the conjugation.

In modern language, put $\sigma(a) = \bar a$. Then $a \mapsto a'$ is a $\sigma$-derivation. The
ring that Ore built is the skew polynomial ring $K[x;\sigma,\delta]$. It is
the subject of Chapter~\ref{ch:generalized-twisted}.

Ore had three special cases in mind, and he lists them~\cite[p.~482]{Ore1933}.
\begin{itemize}
\item Suppose $x$ commutes with the coefficients. Then $\bar a = a$ and
  $a' = 0$. This is the usual polynomial ring.
\item Linear differential equations: here $K$ is commutative and
  $\bar a = a$. So $(ab)' = ab' + a'b$, and ``the derivative has the ordinary
  properties of a derivative''. Ore remarks that this makes it ``possible to
  introduce the notion of a derivative into non-commutative fields''.
\item Difference equations: here $a' = 0$, while $\bar a$ is usually
  different from~$a$.
\end{itemize}

So in 1933 one definition covered three things. It covered the usual
derivations of a field, derivations of a skew field, and the shift operators
of difference equations. The idea of a derivation twisted by an endomorphism
starts here.

Ore studied one derivation at a time. Four years later, Jacobson looked at
all the derivations of an algebra together.
